% load packages
\usepackage[top=0.7in, bottom=2.0in, left=1.25in, right=1in]{geometry}
\usepackage{xcolor}
\usepackage{fontspec}
\usepackage{float}
\usepackage{lscape}
\usepackage{typearea}
\usepackage{tabu, colortbl, multirow, ctable} % crear tablas
\usepackage{tikz}
\usepackage{amsmath}
\usetikzlibrary{arrows,shapes,positioning,shadows,trees}
\usepackage{multicol}
\usepackage{multirow}
\usepackage{array}
\usepackage{ragged2e}
\newcolumntype{P}[1]{>{\RaggedRight\hspace{0pt}}p{#1}}
\newcolumntype{M}[1]{>{\centering\arraybackslash}m{#1}}
\usepackage{pgfplots, pgfplotstable}
\usepackage{caption}
\usepackage{subcaption}
\usepackage{setspace}
\renewcommand{\baselinestretch}{1.5}
\usepackage{graphicx}  % insertar figuras
\usepackage{polyglossia}
\setmainlanguage{spanish}
\usepackage{fontspec}
\defaultfontfeatures{Mapping=tex-text,Scale=MatchLowercase}
%CAMBIAR EL TAMAÑO DE LAS SECCIONES
\usepackage{sectsty}
\sectionfont{\fontsize{14}{15}\selectfont}

%% Let's define some colours
\definecolor{ocre}{RGB}{0,0,146}
\definecolor{color1}{RGB}{255,255,255}
\definecolor{azul}{RGB}{0,0,146}
% DEFINE NUMBER OF NOTE
\newcommand{\NOTA}{No. 144, agosto 2023}	  
%-% Este bloque define los hipervínculos------------------------------------------------------------------------------
\usepackage{hyperref}  % hipervínculos
\hypersetup{
	bookmarks=true,         % show bookmarks bar?
	colorlinks=true,       % false: boxed links; true: colored links
	linkcolor=blue,          % color of internal links
	citecolor=blue,        % color of links to bibliography
	filecolor=magenta,      % color of file links
	urlcolor=blue            % color of external links
}
%  Este bloque diseña el encabezado que aparece en cada página
\usepackage{fancyhdr}
\pagestyle{fancy}
\setlength\headheight{3.5cm}
\rhead{%
	\begin{tabular}{c @{\hspace{0.5cm}} c @{\hspace{2.0cm}}}
	\multirow{5}{*}{\hspace{0.0cm}\includegraphics[width=0.15\textwidth]{Logo.jpg}}&\\
		& \LARGE{\bfseries Notas Económicas Regionales} \\%\cline{2-2}
		&\\
		& \vspace{0.3cm} \hspace{-0.5cm}  \large{Secretaría Ejecutiva del Consejo Monetario Centroamericano} \\
		& \vspace{-0.2cm} \small{\NOTA}
\end{tabular}
}
%-------------------------------------------------------------------------------




%% Use some custom fonts
\setsansfont{Futura}[
   Path=fonts/,
   Scale=1,
   Extension = .ttf,
   UprightFont=*-Light,
   BoldFont=*-Bold,
   ItalicFont=*-Lightitalic,
    ]
\setmainfont{Futura}[
    Path=fonts/,
    Scale=1,
    Extension = .ttf,
    UprightFont=*-Light,
    BoldFont=*-Bold,
    ItalicFont=*-Lightitalic,
    ]

%BIBLIOGRAFÍA
\usepackage[%
bibstyle=authoryear,%
citestyle=authoryear-icomp,%
sorting=nyt,%
backend=bibtex,%
url=false,%
uniquename=false,%
doi=false,%
bibencoding=utf8%
]{biblatex}
\addbibresource{REFERENCIA_3.bib}
\DefineBibliographyStrings{spanish}{
andothers={et~al.\adddot}
}

%  Este bloque redefine el título del documento (portada)
%
\makeatletter
\def\@maketitle{%
	\newpage
	\null
	\vskip 0.5em%
	\begin{center}%
		\let \footnote \thanks
		{\large \@title \par}%
		\vskip 0.7em%
	\end{center}%
	\par
	\vskip 1.5em}
%\fi
\makeatother


%******************FINAL DE LA PLANTILLA************************************************
\usepackage{float}
\usepackage{lscape}

%\onehalfspacing
\title{\Large{Helicóptero fiscal de origen estadounidense en los cielos de Centroamérica y República Dominicana\thanks{Se agradece a Juan Izaguirre, Jorge Madrigal y Odalis Marte, de la Secretaría Ejecutiva del Consejo Monetario Centroamericano, por el apoyo en compartir los datos sobre las exportaciones e importaciones de bienes según destino y los comentarios a una versión preliminar. Las ideas y opiniones en la presente nota no corresponden a las del Consejo Monetario Centroamericano (CMCA) o sus miembros. Los errores u omisiones son responsabilidad exclusiva de su autor.}
}
\\
\lineskip .5em
\begin{tabular}{ccc}
& & \\
\hspace{0.01em}%
     & \normalsize{Luis Ortiz Cevallos\footnote{Economista en la Secretaría Ejecutiva del Consejo Monetario Centroamericano.}} & \\
    \hspace{0.01em}%
    &
		\normalsize{\href{mailto:lortiz@secmca.org}{lortiz@secmca.org}} & \\
\end{tabular}
}  % AJUSTAR TÍTULO


\makeatletter
\pgfdeclareshape{datastore}{
	\inheritsavedanchors[from=rectangle]
	\inheritanchorborder[from=rectangle]
	\inheritanchor[from=rectangle]{center}
	\inheritanchor[from=rectangle]{base}
	\inheritanchor[from=rectangle]{north}
	\inheritanchor[from=rectangle]{north east}
	\inheritanchor[from=rectangle]{east}
	\inheritanchor[from=rectangle]{south east}
	\inheritanchor[from=rectangle]{south}
	\inheritanchor[from=rectangle]{south west}
	\inheritanchor[from=rectangle]{west}
	\inheritanchor[from=rectangle]{north west}
	\backgroundpath{
		%  store lower right in xa/ya and upper right in xb/yb
		\southwest \pgf@xa=\pgf@x \pgf@ya=\pgf@y
		\northeast \pgf@xb=\pgf@x \pgf@yb=\pgf@y
		\pgfpathmoveto{\pgfpoint{\pgf@xa}{\pgf@ya}}
		\pgfpathlineto{\pgfpoint{\pgf@xb}{\pgf@ya}}
		\pgfpathmoveto{\pgfpoint{\pgf@xa}{\pgf@yb}}
		\pgfpathlineto{\pgfpoint{\pgf@xb}{\pgf@yb}}
	}
}
\makeatother




